%=====================================
% Khai báo nhóm Tex (cơ bản)
%=====================================
\usepackage{amsmath,amssymb,mathrsfs,maybemath,xlop,polynom,slashbox}
\usepackage{yhmath} %\let\widering\relax %cần khi sd với font fouriernc
\usepackage{enumerate}
\usepackage{tikz} 
\usepackage{tkz-euclide}
%\usepackage{ex_tkz-euclide}
%\usetkzobj{all}
\usepackage{tikz-3dplot}
\usepackage{tkz-tab}
\usepackage{pifont} %kí hiệu đặc biệt
%\usepackage{bbding}
%\usepackage{array}
%\usepackage{tasks}
%==========
\usetikzlibrary{math,through,calc,intersections,angles,quotes,shapes,shapes.geometric,arrows,patterns,snakes,matrix,chains,arrows.meta,decorations.shapes,decorations.fractals,decorations.markings,shadows}
\usetikzlibrary{positioning,decorations.text,decorations.pathmorphing}% Để uốn cong văn bản 
\usetikzlibrary{shadings,fadings} %ĐỔ BÓNG
\usepackage{pgfplots}
\usepackage{pgfornament}
\usepgfplotslibrary{fillbetween}
\pgfplotsset{compat=1.9}
\usepackage[hidelinks,unicode]{hyperref}
\usepackage{currfile}
\usepackage[outline]{contour} %viền
\usepackage{fontawesome} % Gói kí hiệu
\usepackage{lipsum} %Lấy text
%%---------
%\usepackage{setspace}
%\usepackage{scrextend}
\usepackage{varwidth}
%===========Bảng
\usepackage{longtable,multirow,makecell}
\usepackage{diagbox}
\renewcommand{\tabcolsep}{3mm}
\newcolumntype{C}[1]{>{\centering\arraybackslash}p{#1}}
\newcolumntype{L}[1]{>{\raggedright\arraybackslash}p{#1}}
%-----------Trang vb
\pgfmathsetmacro{\mepphai}{\phai+\leftnote} 
\usepackage[top=\tren cm, bottom=\duoi cm, left=\trai cm, right=\mepphai cm] {geometry}
\usepackage{capt-of}
\usepackage[locale=DE]{siunitx} % cách viết số đo có đơn vị theo chuẩn DE (gần giống VN) 
%--------------Gói trắc nghiệm EX-TEST
\usepackage[loigiai]{ex_test} 
%\usepackage[solcolor]{ex_test}
%--------------------------
%----Lời giải, Hiền thị tên EX; Dấu kết thúc
\font\damEX=ugqb8v at 11pt
%\def\loigiaiEX{\color{\mauLG}\damEX\strut\faCommenting\ Lời giải.}
\def\loigiaiEX{
	\tikz[]{
		\draw (0,0)++(0.5*\textwidth,0) node[inner sep=0pt] {\color{\mauLG}%\fontfamily{qag}\bfseries\strut%\faFolderOpen\ Lời giải.
			\damEX\faCommenting\ Lời giải.};%
	}\vspace*{-2mm}
}
\renewcommand{\nameex}{\fontfamily{pag}\selectfont%
	\damEX\color{\mauEX} Câu}
\def\mauVuong{cyan}
\def\qedEX{\color{\mauVuong}\ensuremath{\square}}
%---------- Khai báo viết tắt, in đáp án
\newcommand{\hoac}[1]{ %hệ hoặc
	\left[\begin{aligned}#1\end{aligned}\right.}
\newcommand{\heva}[1]{ %hệ và
	\left\{\begin{aligned}#1\end{aligned}\right.}
%--In đáp án
\newcommand{\indapan}[2]{
	\addcontentsline{toc}{subsection}{\sf Bảng đáp án} % đưa MT vào mục lục
	\par\vspace*{5mm}
	\begin{tikzpicture}%
		\draw (0,0)++(0.5*\textwidth,0) node[thick,scale=1,fill=\mauEX!2,draw=\maufoot,minimum width=3.5cm,minimum height=0.1cm,rounded corners=2mm] {\damEX\color{\mauEX} BẢNG ĐÁP ÁN};
	\end{tikzpicture}%
	%		\end{center}
\vspace*{-4mm}
\inputansbox{#1}{#2}
}
%--In đáp án True
\newcommand{\indapanT}[3]{
\addcontentsline{toc}{subsection}{\sf Bảng đáp án 1} % đưa MT vào mục lục
\par%\vspace*{5mm}
\begin{tikzpicture}%
	\draw (0,0)++(0.5*\textwidth,0) node[thick,scale=1,fill=\mauEX!2,draw=\maufoot,minimum width=3.5cm,minimum height=0.1cm,rounded corners=2mm] {\damEX\color{\mauEX} #3};
\end{tikzpicture}%
%		\end{center}
\vspace*{-4mm}
\inputansbox{#1}{#2}
}
%--In đáp án TF
\newcommand{\indapanTF}[3]{
\addcontentsline{toc}{subsection}{\sf Bảng đáp án 2} % đưa MT vào mục lục
\par%\vspace*{5mm}
\begin{tikzpicture}%
\draw (0,0)++(0.5*\textwidth,0) node[thick,scale=1,fill=\mauEX!2,draw=\maufoot,minimum width=3.5cm,minimum height=0.1cm,rounded corners=2mm] {\damEX\color{\mauEX} #3};
\end{tikzpicture}%
%		\end{center}
\vspace*{-4mm}
\inputansboxTF{#1}{#2}
}
%--In đáp án SA
\newcommand{\indapanSA}[3]{
\addcontentsline{toc}{subsection}{\sf Bảng đáp án 3} % đưa MT vào mục lục
\par%\vspace*{5mm}
\begin{tikzpicture}%
\draw (0,0)++(0.5*\textwidth,0) node[thick,scale=1,fill=\mauEX!2,draw=\maufoot,minimum width=3.5cm,minimum height=0.1cm,rounded corners=2mm] {\damEX\color{\mauEX} #3};
\end{tikzpicture}%
%		\end{center}
\vspace*{-4mm}
\inputansboxSA{#1}{#2}
}
%----------
\usepackage{esvect}
\def\vec{\vv} %vecto
\def\overrightarrow{\vv}
%Lệnh song song
\DeclareSymbolFont{symbolsC}{U}{txsyc}{m}{n}
\DeclareMathSymbol{\varparallel}{\mathrel}{symbolsC}{9}
\DeclareMathSymbol{\parallel}{\mathrel}{symbolsC}{9}
%--------------------------
% HEADER AND FOOTER STYLING
%--------------------------
%Tạo background
\def\hienLOGO{
\usepackage{background}
\backgroundsetup{%
scale=1.25,%%
angle=0,%%
opacity=0.1,%%
contents={
\ifnum\the\value{page}>0%
\begin{tikzpicture}[remember picture,overlay,scale=1.25]
\node at (current page.center){\includegraphics[width=13cm]{logo/logonen.jpg}};
\end{tikzpicture}
\else
\fi
}
}	
}
%----
\def\anLOGO{}
%--------------------------
%\def\trenletrai{\rightmark}
%\def\trenchanphai{\leftmark}
\newcommand{\myfancyhead}{% trên
\begin{tikzpicture}[remember picture,overlay]
\checkoddpage\ifoddpage %nếu trang lẻ
\def\kctrai{\trai} \def\kcphai{\phai}
\pgfmathsetmacro{\kctren}{-\tren+0.5*\topset} 
\path ([xshift=\kctrai cm, yshift=\kctren cm]current page.north west) coordinate (AA)
++(\textwidth+\leftnote cm,0)coordinate (BB); 
%---- màu bên trái
\foreach \ii/\mau in {1.35/\maufoot!30,1/\maufoot!50,0.5/\maufoot!70,-0.1/\maufoot}{
\fill[\mau,draw=white,line width=2pt] ([xshift=-1pt,yshift=1pt]AA)--++(0.75 cm+\ii cm+3pt,0)--++(0.3,0.5cm+1pt)--++(-0.75 cm -\ii cm-3pt-0.3cm,0)--cycle;
}
%---- chữ
\node[anchor=south west,text=\maufoot,inner sep=0pt] at ([xshift=0.75 cm+1.35cm+10pt,yshift=1pt]AA){\fontfamily{qag}\fontsize{10pt}{12pt}\selectfont\nouppercase{\trenletrai}};
%-----đường kẻ
\draw[\maukefoot, line width=1.5pt] (AA) -- (BB);
%-----bên phải
\node[rectangle, fill=\maufoot, text=white, anchor=south east,inner sep=3pt, rounded corners=4pt] (sotrang) at ([xshift=0cm,yshift=1.5pt]BB){\fontfamily{qag}\fontsize{10pt}{12pt}\selectfont\bfseries\thepage};
\node[text=\maufoot, anchor=east,inner sep=0pt] at ([xshift=-3pt,yshift=-1pt]sotrang.west){\fontfamily{qag}\fontsize{10pt}{18pt}\selectfont Trang};
%-------------------
\else % trang chẵn
\def\kctrai{\phai} \def\kcphai{\trai}
\pgfmathsetmacro{\kctren}{-\tren+0.5*\topset} 
\path ([xshift=\kctrai cm, yshift=\kctren cm]current page.north west) coordinate (AA)
++(\textwidth+\leftnote cm,0)coordinate (BB); 
%---- bên phải
\foreach \ii/\mau in {1.35/\maufoot!30,1/\maufoot!50,0.5/\maufoot!70,-0.1/\maufoot}{
\fill[\mau,draw=white,line width=2pt] ([xshift=1pt,yshift=1pt]BB)--++(-0.75 cm-\ii cm-3pt,0)
--++(-0.3,0.5cm+1pt)--++(0.75 cm +\ii cm+3pt+0.3 cm,0)--cycle
;
}
\node[anchor=south east,text=\maufoot,inner sep=0pt] at ([xshift=-0.75 cm-1.35cm-10pt]BB){\fontfamily{qag}\fontsize{10pt}{12pt}\selectfont\nouppercase{\trenchanphai}};
%-----đường kẻ
\draw[\maukefoot, line width=1.5pt] (AA) --(BB);
\node[text=\maufoot, anchor=west,inner sep=0pt] (trang) at ([yshift=6pt]AA){\fontfamily{qag}\fontsize{10pt}{12pt}\selectfont Trang};
\node[rectangle, fill=\maufoot, text=white, anchor=west,inner sep=3pt, rounded corners=4pt] at ([xshift=3pt,yshift=1.5pt]trang.east){\fontfamily{qag}\fontsize{10pt}{18pt}\selectfont\bfseries\thepage};
\fi
%------
%=====================Dòng chấm note
\path ([yshift=-\tren cm+\topset cm]current page.north west) coordinate (AAA)
++(\paperwidth,0)coordinate (BBB);
\checkoddpage\ifoddpage %nếu trang lẻ
%----- Kẻ đứng
\draw[\maukefoot] ([xshift=-\mepphai cm+2.5mm,yshift=-3mm]BBB)--([yshift=\duoi cm-0.75*\botset cm+0mm,xshift=-\mepphai cm+2.5mm]current page.south east);
%--		
\path ([yshift=-\tren cm+0.5*\topset cm-0.35cm,xshift=-\phai cm-0.5*\leftnote cm+2.5mm]current page.north east) coordinate (DDD); 
\begin{scope}
\clip ([yshift=-\tren cm+0.5*\topset cm-1pt,xshift=-\phai cm]current page.north east) rectangle ([yshift=\duoi cm-0.5*\botset cm,xshift=-\mepphai cm+5mm]current page.south east);% cắt chấm
\node[inner sep =0pt,scale=1,anchor=north] at ([yshift=0cm,xshift=0pt]DDD) {
\parbox{\leftnote cm}{\centering
	\dotlineEXhead{60}
}
};
\end{scope}
\else %chẵn
%----- Kẻ đứng
\draw[\maukefoot] ([xshift=\mepphai cm-2.5mm,yshift=-3mm]AAA)--([yshift=\duoi cm-0.75*\botset cm+0mm,xshift=\mepphai cm-2.5mm]current page.south west);
%--		
\path ([yshift=-\tren cm+0.5*\topset cm-0.35cm,xshift=\phai cm+0.5*\leftnote cm-2.5mm]current page.north west) coordinate (DDD); 
\begin{scope}
\clip ([yshift=-\tren cm+0.5*\topset cm-1pt,xshift=\phai cm]current page.north west) rectangle ([yshift=\duoi cm-0.5*\botset cm,xshift=\mepphai cm-5mm]current page.south west);% cắt chấm
\node[inner sep =0pt,scale=1,anchor=north] at ([yshift=0cm,xshift=0pt]DDD) {
\parbox{\leftnote cm}{\centering
	\dotlineEXhead{60}
}
};
\end{scope}					
\fi
%--note dưới
\node[inner sep =6pt, text=white,scale=1,anchor=north,fill=\maufoot] (noteduoi) at ([yshift=2pt]DDD) {
\parbox{\leftnote cm-5mm-12pt}{\fontfamily{qag}\fontsize{11}{1}\selectfont\bfseries\centering
QUICK NOTE
}
};
\draw[\maufoot, line width=0.4pt] ([yshift=-2pt]noteduoi.south west)--([yshift=-2pt]noteduoi.south east);
\end{tikzpicture}%
}
% trên mục lục
\newcommand{\headmucluc}{%
\boldmath
\begin{tikzpicture}[remember picture,overlay,>=stealth]
\checkoddpage\ifoddpage %nếu trang lẻ
\def\kctrai{\trai} \def\kcphai{\phai}
\pgfmathsetmacro{\kctren}{-\tren+0.5*\topset} 
\path ([xshift=\kctrai cm, yshift=\kctren cm]current page.north west) coordinate (AA)
++(\textwidth,0)coordinate (BB); 
%---- màu bên trái
\foreach \ii/\mau in {1.35/\maufoot!30,1/\maufoot!50,0.5/\maufoot!70,-0.1/\maufoot}{
\fill[\mau,draw=white,line width=2pt] ([xshift=-1pt,yshift=1pt]AA)--++(0.75 cm+\ii cm+3pt,0)--++(0.3,0.5cm+1pt)--++(-0.75 cm -\ii cm-3pt-0.3cm,0)--cycle;
}
%---- chữ
\node[anchor=south west,text=\maufoot,inner sep=0pt] at ([xshift=0.75 cm+1.35cm+10pt,yshift=1pt]AA){\fontfamily{qag}\fontsize{10pt}{12pt}\selectfont\nouppercase{\trenletrai}};
%-----đường kẻ
\draw[\maukefoot, line width=1.5pt] (AA) -- (BB);
%-----bên phải
\node[rectangle, fill=\maufoot, text=white, anchor=south east,inner sep=3pt, rounded corners=4pt] (sotrang) at ([xshift=0cm,yshift=1.5pt]BB){\fontfamily{qag}\fontsize{10pt}{12pt}\selectfont\bfseries\thepage};
\node[text=\maufoot, anchor=east,inner sep=0pt] at ([xshift=-3pt,yshift=-1pt]sotrang.west){\fontfamily{qag}\fontsize{10pt}{18pt}\selectfont Trang};
%-------------------
\else % trang chẵn
\def\kctrai{\phai} \def\kcphai{\trai}
\pgfmathsetmacro{\kctren}{-\tren+0.5*\topset} 
\path ([xshift=\kctrai cm, yshift=\kctren cm]current page.north west) coordinate (AA)
++(\textwidth,0)coordinate (BB); 
%---- bên phải
\foreach \ii/\mau in {1.35/\maufoot!30,1/\maufoot!50,0.5/\maufoot!70,-0.1/\maufoot}{
\fill[\mau,draw=white,line width=2pt] ([xshift=1pt,yshift=1pt]BB)--++(-0.75 cm-\ii cm-3pt,0)
--++(-0.3,0.5cm+1pt)--++(0.75 cm +\ii cm+3pt+0.3 cm,0)--cycle
;
}
\node[anchor=south east,text=\maufoot,inner sep=0pt] at ([xshift=-0.75 cm-1.35cm-10pt]BB){\fontfamily{qag}\fontsize{10pt}{12pt}\selectfont\nouppercase{\trenchanphai}};
%-----đường kẻ
\draw[\maukefoot, line width=1.5pt] (AA) --(BB);
\node[text=\maufoot, anchor=west,inner sep=0pt] (trang) at ([yshift=6pt]AA){\fontfamily{qag}\fontsize{10pt}{12pt}\selectfont Trang};
\node[rectangle, fill=\maufoot, text=white, anchor=west,inner sep=3pt, rounded corners=4pt] at ([xshift=3pt,yshift=1.5pt]trang.east){\fontfamily{qag}\fontsize{10pt}{18pt}\selectfont\bfseries\thepage};
\fi
\end{tikzpicture}%
}
%===========================
\newcommand{\myfancyfoot}{% dưới
\begin{tikzpicture}[remember picture,overlay]
\checkoddpage\ifoddpage %nếu trang lẻ
\def\kctrai{\trai} \def\kcphai{\phai}
\else
\def\kctrai{\phai} \def\kcphai{\trai}
\fi
\pgfmathsetmacro{\kcduoi}{\duoi-\botset} 
\path ([yshift=\kcduoi cm,xshift=\kctrai cm]current page.south west) coordinate (AA)
([yshift=\kcduoi cm,xshift=-\kcphai cm]current page.south east)coordinate (BB); 
%-----đường kẻ
\draw[\maukefoot,line width=1.5pt] (AA) --(BB);
%-----nền
\fill[left color=\maukefoot!20, right color=\maukefoot!20, middle color=white]([yshift=-1mm]AA) rectangle ([yshift=-7mm]BB);
%-----chính giữa
%		\node[anchor=north,text=\maufoot,inner sep=0pt,text width=0.33*\textwidth+\leftnote cm,align=center] (dc) at ([yshift=-4pt]$(AA)!0.5!(BB)$){\fontfamily{qag}\fontsize{9pt}{1pt}\selectfont\bfseries
%			 \noidungduoigiua%
%		 };
\draw[line width=0.5pt,fill=white,draw=\maukefoot,path picture={\node at ([yshift=-4mm]$(AA)!0.5!(BB)$) {\includegraphics[width=11mm]{logo/logo.jpg}};}] ([yshift=-4mm]$(AA)!0.5!(BB)$) circle [radius=7mm]; 
%-----bên trái
\node[anchor=north west,text=\maufoot,inner sep=0pt,text width=0.5*\textwidth,align=left] (sdt) at ([yshift=-4pt]AA){\fontfamily{qag}\fontsize{10pt}{1pt}\selectfont
\noidungduoitrai%
};
%------bên phải
\node[anchor=north east,text=\maufoot,inner sep=0pt,text width=0.5*\textwidth,align=right] (face) at ([yshift=-4pt]BB){\fontfamily{qag}\fontsize{10pt}{1pt}\selectfont
\noidungduoiphai%
};
\end{tikzpicture}%
}
%----------------------
\usepackage{changepage}
\strictpagecheck
\usepackage{lastpage}
\usepackage{fancyhdr,lastpage}
\pagestyle{fancy}
\fancyhf{}
\def\headchuong{\myfancyhead}
\fancypagestyle{plain}{
\fancyhead[LO,RE]{\headchuong}
\fancyfoot[LO,RE]{\myfancyfoot}
}
\fancyhead[LO,RE]{\myfancyhead}
\fancyfoot[LO,RE]{\myfancyfoot}
\renewcommand{\footrulewidth}{0pt}
\renewcommand{\headrulewidth}{0pt}
%--------------4.2
\usepackage[most]{tcolorbox}
\colorlet{tcbcol@back}{tcbcolback}
\colorlet{tcbcol@frame}{tcbcolframe}
%-------------- Khung nội dung 2 tham số
\newtcolorbox{noidung}[1]{
breakable, enhanced jigsaw,	frame hidden,opacityback=0,
%	enhanced jigsaw,breakable,pad at break*=1mm,
%	colback=pink!40,
%	%opacityback=0, %ko nền
%	colframe=pink!40,
coltitle=\maund,
left=14mm,right=3mm,top=2mm,bottom=2mm,
%--
before skip=2mm,
after skip=3mm,
boxrule=0.5pt, %đường viền dày
%toprule=1pt,
%bottomrule=1pt,
%leftrule=2.5pt,
%rightrule=0pt,
%	sharp corners,
%	rounded corners=southwest, 
%	rounded corners=northeast, 
%	rounded corners=southeast,
%	arc=7mm,
underlay unbroken and first={
%		\draw[draw=cyan,line width=1pt] ([xshift=6.5mm]title.north west) -- (title.north east);
%--khung nội dung
\fill[fill=\maund!4,drop shadow=\maund!30,opacity=0.5] 
([xshift=13mm]interior.north west) arc (0:-90:8mm) 
--([xshift=5mm]interior.south west)
--([xshift=-8mm]interior.south east) arc (-90:0:8mm) 
--(interior.north east)
--cycle;
\draw[\maund,line width=1pt] 
(interior.north east) coordinate (mot)		
--([xshift=13mm]interior.north west) coordinate (hai) arc (0:-90:8mm) coordinate (ba) 
--([xshift=5mm]interior.south west) coordinate (bon)
;
\fill[\maund] 
(mot) circle (2pt)
(hai) circle (2pt)
(ba) circle (2pt)
(bon) circle (2pt)
;
},
underlay last={
%--khung nội dung
\fill[fill=\maund!2,drop shadow=gray] 
([xshift=5mm]interior.north west)
--([xshift=5mm]interior.south west)
--([xshift=-8mm]interior.south east) arc (-90:0:8mm) 
--(interior.north east)
--cycle;
\draw[\maund,line width=1pt] 
(interior.north east) coordinate (mot)		
--([xshift=5mm]interior.north west) coordinate (hai) 
--([xshift=5mm]interior.south west) coordinate (bon)
;
\fill[\maund] 
(mot) circle (2pt)
(hai) circle (2pt)
(bon) circle (2pt)
;
},
fonttitle=\fontfamily{qag}\bfseries\selectfont,
title={#1},
%	attach boxed title to top left={xshift=1cm,yshift=-1mm},
%	boxed title style={empty},
%	underlay boxed title={
underlay unbroken and first={
\fill[\maund!70] ([xshift=5mm]interior.north west) node[white,scale=1.5]{
%\faBell
\faEdit
} circle (6.5mm) ;
\draw[white] ([xshift=5mm]interior.north west) circle (6mm) ;
},
%----
%fontupper=\sffamily % chữ nghiêng phần nội dung
}
%-------------- Khung 
\newtcolorbox[auto counter]{khung4}[1]{enhanced jigsaw, breakable,
before skip=3mm,after skip=3mm,
left=1mm,right=1mm,top=2mm,bottom=1mm,
colframe=\maukhung4,colback=\maukhung4!2,colbacktitle=\maukhung4!6,coltitle=\maukhung4!90!black,
opacityback=0.5,
boxrule=0.4mm,
attach boxed title to top center=
{yshift=-0.1mm-\tcboxedtitleheight/2,yshifttext=2mm-\tcboxedtitleheight/2},
varwidth boxed title*=-3cm,
boxed title style={boxrule=0.3mm,
frame code={ \path[tcb fill frame] ([xshift=-4mm]frame.west)
-- (frame.north west) -- (frame.north east) -- ([xshift=4mm]frame.east)
-- (frame.south east) -- (frame.south west) -- cycle; },
interior code={ \path[tcb fill interior] ([xshift=-2mm]interior.west)
-- (interior.north west) -- (interior.north east)
-- ([xshift=2mm]interior.east) -- (interior.south east) -- (interior.south west)
-- cycle;} 
},
fonttitle=\fontfamily{qag}\fontsize{10}{0}
\bfseries,
fontupper=\sffamily,
title={#1}
}
%---------------Dạng toán
\newcounter{dang}\setcounter{dang}{0}
\renewcommand{\thedang}{\arabic{dang}}
%%%%%%%% KIỂU TUỲ CHỌN KHUNG
\newcommand{\TuyChonDang}[1]{
\def\chondang{#1}
\def\dmot{1}
\def\dhai{2}
\ifx\chondang\dmot%
%---Dạng 1
\newtcolorbox{dang}[1]{
fonttitle=\fontfamily{qag}\bfseries,%fontupper=\itshape,
breakable, 
enhanced jigsaw, %nâng cao
%	frame hidden, opacityback=0,
colbacktitle=\nensodang,
coltitle=white,
colframe=\maudang,
colback=\nendang,
opacityback=0.5,
sharp corners,
before skip=3mm,after skip=3mm,
left=2mm,right=2mm,top=2mm,bottom=2mm,
%	boxrule=1pt,
%--cong góc
%	rounded corners=southwest, 
%	rounded corners=northeast, 
%	rounded corners=northwest,
rounded corners=southeast,
arc is angular,
arc=3mm,
underlay unbroken and first={
\def\kctrai{1.5}
\def\daikhung{1}
%--nền
\fill[\maudang] 
(title.north west) rectangle ([xshift=\kctrai cm+\daikhung cm+2.5mm]title.south west)
;
%--Nền số
\fill[\nensodang] 
([xshift=\kctrai cm,yshift=0.12cm+0.1cm]title.north west) coordinate (bai1)%trên trái
--++(\daikhung,0) coordinate (bai2)%trên phải
[rounded corners=6pt]--++(0,-0.6) coordinate (bai3)%dưới phải
--++(-\daikhung,0) coordinate (bai4)%dưới trái
[rounded corners=0pt]--cycle
;
%--Bóng góc trái trên của khung
\fill[gray!50] (bai1)--++(-0.2,-0.15)--++(0.2,0)--cycle;
%--Bài và số bài
\node[circle,draw=white,fill=\sodang, line width=1pt, inner sep=3pt,text=white] (SOBAI) at ([yshift=1pt]$(bai4)!0.5!(bai3)$) {\fontsize{12.5pt}{1pt}\fontfamily{put}\selectfont\bfseries \thedang};
\draw ([yshift=-1.5pt]SOBAI.west) node[left,white,scale=0.85] {\fontsize{13pt}{1pt}\fontfamily{qag}\selectfont\bfseries DẠNG};	
},
title={%\faFolderOpen\ Dạng~\stepcounter{dang}\thedang.\ 
\hspace*{2.75cm}\parbox{\textwidth-2.75cm}{##1}\stepcounter{dang}}%
\addcontentsline{toc}{subsection}{\it\sffamily \faFolderOpen\ Dạng~\thedang. ##1}
}
\else
%--------
\newtcolorbox{dang}[1]{
breakable, 
%	blanker, %enhanced, 
%--
enhanced jigsaw,
frame hidden,opacityback=0,
%--
before skip=6mm,after skip=3mm,
left=2mm,right=2mm,top=4mm,bottom=2mm,
%colback=green!10,colframe=cyan,boxrule=0.3mm,
underlay={
%--khung nội dung
\fill[draw=\maudang, top color=\nendang, bottom color=\nendang, middle color=white, rounded corners=0mm, line width=1pt,drop shadow=\nensodang!50] ([yshift=-2mm]interior.north west) rectangle (interior.south east);
},
underlay unbroken and first={
%--khung tiêu đề
\fill[\maudang,rounded corners=0mm,draw=\maudang,drop shadow=\nensodang] ([yshift=0mm,xshift=0cm]title.north west) rectangle ([yshift=0mm,xshift=0cm]title.south east);
\draw[white, rounded corners=0mm] ([yshift=-.3mm,xshift=0.3mm]title.north west) rectangle ([yshift=0.3mm,xshift=-0.3mm]title.south east);
%--khung dạng
%bóng trái
\fill[gray] ([xshift=0.75cm,yshift=2mm]title.north west)--++(-2mm,-2mm)--++(2mm,0)--cycle;
%khung
\fill[\nensodang,drop shadow] 
([xshift=0.75cm,yshift=2mm]title.north west)
--++(1.2,0)
[rounded corners =4pt]--++(0,-0.9)
--++(-1.2,0)
[rounded corners =0pt]--cycle;
\node[white, below] (dang) at ([xshift=1.35cm,yshift=2.5mm]title.north west) {\fontfamily{qag}\fontsize{6pt}{0pt}\bfseries DẠNG};
\draw ([yshift=-1pt]dang) node[below,white] {\fontfamily{qag}\fontsize{15pt}{0pt}\bfseries \thedang};
},
fonttitle=\fontfamily{qag}\fontsize{10pt}{0pt}\selectfont\bfseries, %fontupper=\itshape,
halign title=center, %canh giữa
title={%\faFolderOpen\ Dạng~\stepcounter{dang}\thedang.\ 
\hspace*{2cm}\parbox{\textwidth-1.8cm}{\centering ##1}\stepcounter{dang}}%
\addcontentsline{toc}{subsection}{\it\sffamily \faFolderOpen\ Dạng~\thedang. ##1}
}
\fi
}
%---------------------------------------------------------------
% ĐỊNH NGHĨA SECTION. SUBSECTION, SUBSUBSECTION ... THEO Ý RIÊNG
%---------------------------------------------------------------
\usepackage[explicit]{titlesec} % để gọi #1
\usepackage{titledot} % gói lệnh chứa cả titlesec và titletoc
%=====================================
\setcounter{secnumdepth}{4} %độ sâu
\renewcommand{\thechapter}{\arabic{chapter}}
\renewcommand{\thesection}{\arabic{section}}
\renewcommand{\thesubsection}{\Alph{subsection}}
\renewcommand{\thesubsubsection}{\arabic{subsubsection}}
%--------------Tròn
\newcommand{\tron}[1]{
\begin{tikzpicture}[baseline=(A.base)]%
\node[circle,draw=\mauSUBSEC,line width=0.5pt,fill=white,inner sep=2pt,outer sep=1pt] (A) {\color{white} #1};
\node[circle,draw=none,fill=\mauSUBSEC,inner sep=1pt,outer sep=1pt] (A) {\color{white} #1};
\end{tikzpicture}%
}
%-- Số chương
%--------------Tròn
\newcommand{\sochuong}[1]{
\node[fill=\maunenkhoi,inner sep=1pt,text=white,anchor=center,rounded corners=6pt,drop shadow={shadow yshift=0pt,shadow xshift=1pt,black!90},drop shadow={shadow yshift=0pt,shadow xshift=-1pt,black!90}] at ($(Ctrai)!0.5!(Cphai)$){%
\fontfamily{put}\bfseries\fontsize{11}{11}\selectfont #1%
};
}
%================= Đn chương
%\def\khoi{12}
%\def\tieudetrenchuong{ÔN THI THPT QUỐC GIA 2021-2022}
%\def\maunenchuong{cyan!80!black}
%\def\maunenkhoi{magenta}
\def\cochuchuong{16pt}
\font\damTC=ugqb8v at 16pt
\font\damC=ugqb8v at \cochuchuong
\def\kchaiben{2.75}
\pgfmathsetmacro{\Rchuong}{1.15*\kchaiben/3}
\pgfmathsetmacro{\kcphu}{2*\kchaiben+2*0.85}
\pgfmathsetmacro{\dailogo}{1.5*\Rchuong}
\def\cao{2}
\titlespacing{\chapter}{0cm}{0.5cm}{1.5cm}[0cm] %1: , 2: Trên, 3: dưới
\titleformat{\chapter}
{\fontsize{16pt}{16pt}\fontfamily{qag}\selectfont\bfseries\color{white}} %định dạng chung
{} %đánh số
{-0.4cm}
{
\parbox{\textwidth-\kcphu cm-10pt}%
{\damC\centering\MakeUppercase{#1}} %neo vitri
\begin{tikzpicture}[remember picture,overlay,line join=round,line cap=round] 
%%------------------------------------
\checkoddpage\ifoddpage %nếu trang lẻ
\def\sokc{\mepphai}
\else
\def\sokc{\trai}
\fi
%==============
%-- nền dưới
\node[anchor=north east,inner sep=0pt,fill=\maunenchuong, draw=blue, text width=\textwidth,minimum height=3cm,rounded corners=10pt,drop shadow={shadow yshift=-3pt,shadow xshift=0pt,black!50}]  (tam) at ([xshift=-\sokc cm,yshift=-1.25*\tren cm]current page.north east){};
\draw[\mauchuchuong, line width=1pt,rounded corners=8
pt] ([yshift=-3pt, xshift=3pt]tam.north west) rectangle ([yshift=3pt, xshift=-3pt]tam.south east);
%-- nền trên
\fill[\maunenchuong, draw=\mauchuchuong!90!black, line width =1.5pt,drop shadow={shadow yshift=0pt,shadow xshift=3pt,black!90},drop shadow={shadow yshift=0pt,shadow xshift=-3pt,black!90}]%
([xshift=\kchaiben cm]tam.west) coordinate (Ctrai)--++(0.85,\cao)--++(\textwidth-2*\kchaiben cm-2*0.85cm,0)--([xshift=-\kchaiben cm]tam.east) coordinate (Cphai) --++(-0.85,-\cao)--++(-\textwidth+2*\kchaiben cm+2*0.85cm,0)--cycle;
\draw[\mauchuchuong,line width=0.5pt,dashed] (Ctrai)--(Cphai);
%--- lớp
\fill[\maunenkhoi,drop shadow={shadow yshift=0pt,shadow xshift=1pt,black!90}] ([xshift=0.5*\kchaiben cm]tam.west) circle (\Rchuong);
\node[scale=2.2,circle,text=\maunenchuong,inner sep=4pt] at ([xshift=0.5*\kchaiben cm]tam.west){%
\contour{white}{%
	\khoi%
}%
};
%--- logo chuong
\draw[fill=white,draw=none,drop shadow={shadow yshift=0pt,shadow xshift=-1pt,black!90},path picture={\node at ([xshift=-0mm-0.5*\kchaiben cm,yshift=0mm]tam.east){\includegraphics[width=\dailogo cm]{logo/logo.jpg}};}] ([xshift=-0.5*\kchaiben cm]tam.east) circle [radius=\Rchuong cm]; 
%==========nội dung chương
\node[inner sep=5pt,text=\mauchuchuong,anchor=center] at ([yshift=-0.5*\cao cm]$(Ctrai)!0.5!(Cphai)$){%
\parbox{\textwidth-\kcphu cm-10pt}%
{\damC\centering\MakeUppercase{#1}}
};
%-- Tiêu đề trên chương
\node[inner sep=5pt,text=white,anchor=center] at ([yshift=0.5*\cao cm+3pt]$(Ctrai)!0.5!(Cphai)$){%
\parbox{\textwidth-\kcphu cm-10pt}%
{\damTC\centering\tieudetrenchuong}
};
%--	số chương
\sochuong{\hspace*{2mm}\chaptername ~\thechapter\hspace*{2mm}}
\end{tikzpicture}
}
[\vspace{0.5cm}]
%============================Mục lục - Chapter*
%\makeatletter
\titleformat{name=\chapter,numberless}[display]
{\fontsize{21pt}{18pt}\fontfamily{qag}\selectfont\bfseries\color{\maunenchuong}} %định dạng chung
{}
{-3em}
{%
\begin{tikzpicture}
%-----Nội dung
\node[inner sep=0pt,right,scale=1.65] (ndchuong) at (0,0){\usefont{T5}{HLButlong}{m}{n} %\MakeUppercase
{#1}};
%-----Đường kẻ ngang
\begin{scope}
\clip (0,-0.75) rectangle +(\textwidth,1.5);
\draw[\maunenchuong,line width=2pt] (ndchuong.south east)++(10pt,8pt) --++(\linewidth,0);
\end{scope}
\end{tikzpicture}
}
[%
\vspace{-2cm}%
\thispagestyle{empty}%
]%
%================== ĐN Part
\font\damP=ugqb8v at 22pt
\titlespacing{\part}{0cm}{0cm}{0cm}[0cm]
\titleformat{\part}[display]
{\bfseries}
{}
{0pt}
{
\begin{tikzpicture}[remember picture, overlay]
\checkoddpage\ifoddpage %nếu trang lẻ
\def\kcphu{\phai}
\else
\def\kcphu{\trai}
\fi
\fill[\mauphan!5] (current page.north west) rectangle (current page.south east);
%----Tên phần
\node[inner sep=0pt,scale=2,left] (tenphan) at ([xshift=-\kcphu cm,yshift=-7cm]current page.north east) {\color{\mauphan!70!black}\fontfamily{qag}\bfseries\fontsize{15pt}{1pt}\selectfont\MakeUppercase{\partname}\fontsize{100pt}{1pt}\selectfont \thepart};
%----Nền màu trái
\foreach \ii/\mau in {0/\mauphan!30,0.5/\mauphan!70,1.2/\mauphan}{
\fill[\mau,draw=\mauphan!5,line width=4pt] 
([xshift=-4mm-\ii cm,yshift=1cm]tenphan.south west)
--++(0,-1.2)--++(-\paperwidth,0)--++(0,1.2)--cycle;
}
%-----Nội dung phần
\node[inner sep=0pt,below left,scale=4] at ([xshift=0cm,yshift=-0.5cm]tenphan.south east) 
{
\parbox{5.5cm}{
	\fontsize{20pt}{20pt}\fontfamily{qag}\selectfont\bfseries\color{\mauphan}\raggedleft
	\damP\MakeUppercase{#1}
}
};
\end{tikzpicture}
}
[
\vspace{2cm}
\thispagestyle{empty}
]
%--------Đn Section---------------------------
\newcounter{stt}\setcounter{stt}{0}
\renewcommand{\thestt}{\arabic{stt}}
%\def\mautraibai{magenta}
%\def\mauphaibai{cyan}
\def\tenlop{Lớp vật lý Cô Thảo - Thầy Sang}
\def\slogan{NĂM HỌC 2024 - 2025}
%
\titlespacing*{\section}{0cm}{0cm}{0cm}[0cm]
\titleformat
{\section}
{\color{white}\fontfamily{qag}\fontsize{13pt}{13pt}\selectfont\bfseries}
{}
{0cm}
{%	
\setcounter{ex}{0}\setcounter{EX}{0}%
\setcounter{vd}{0}\setcounter{VD}{0}%
\setcounter{bt}{0}\setcounter{BT}{0}%
\setcounter{dang}{0}%
\setcounter{stt}{0}%
\begin{tikzpicture}%[line join=round,line cap=round]
\def\dddd{4}
\def\ndbai{\centering\S\thesection. ~\MakeUppercase{#1}}
%-- Nội dung chính --	
\node[inner sep=5pt,text=white,fill=\mauphaibai,text width=\textwidth-\dddd cm,align=center] (secname) at (0,0){\ndbai};
%-- số section
\fill[\mautraibai] 
([xshift=-\dddd cm+11pt]secname.north west) coordinate (bai1)%trên trái
rectangle ([xshift=-1pt]secname.south west)coordinate (bai2)
;
%-- số bài
\node[inner sep=0pt,text=white] at ($(bai1)!0.5!(bai2)$) {\fontsize{11pt}{1pt}\fontfamily{qag}\selectfont\bfseries K\khoi ~--~\MakeUppercase{\chaptername ~\thechapter}};
%%-- tên lớp và slogan
%\draw[\mautraibai, line width=2pt] ([yshift=0.4cm,xshift=-\dddd cm+11pt]secname.north west) coordinate (loptoan) --++(\textwidth-3pt,0);
%\node[anchor=west, text=\mautraibai, fill=white, inner sep=0pt,scale=1.1] at ([xshift=-2pt]loptoan){
%\usefont{T5}{HLBrush2BK}{m}{n}\tenlop: ~%
%};
%\node[anchor=center, text=\mautraibai, fill=white, inner sep=3pt] at  ([yshift=0.4cm,xshift=-0.5*\textwidth+0.5*\dddd cm]secname.north east) {%
%\fontfamily{qag}\fontsize{12}{12}\selectfont\slogan%
%};
%-- HS và lớp
%\draw[\mautraibai, line width=0.5pt,dotted] ([yshift=-0.45cm,xshift=-\dddd cm+11pt]secname.south west) coordinate (HS) --++(\textwidth-1pt,0);
%\node[anchor=west, text=\mautraibai, fill=white, inner sep=0pt,scale=1.1] at ([xshift=-2pt,yshift=2pt]HS){
%\usefont{T5}{HLBrush2BK}{m}{n}Họ và tên học sinh: ~%
%};
%\node[anchor=center, text=\mautraibai, fill=white, inner sep=3pt] at  ([yshift=-0.45cm+2pt,xshift=-0.4*\textwidth+0.5*\dddd cm]secname.south east) {%
%\usefont{T5}{HLBrush2BK}{m}{n}Lớp:~%
%};
\end{tikzpicture}}
[]
%-------Đn subsection---------------------------
%\titlespacing{\subsection}{0cm}{0cm}{0cm}[0cm]
%\titleformat{\subsection}
%{\normalfont\fontsize{15pt}{20pt}\fontfamily{put}\selectfont\bfseries\color{\mausubsec}}
%{\hspace*{-2mm}\dnsub{\thesubsection}}
%{2mm}
%{#1}
%[]
%%%===============================================
\newcommand{\dnsub}[1]{
\begin{tikzpicture}[baseline=(muccon.base),line join=round,line cap=round]%
\node[rectangle,inner sep=7pt,text=\mausubsec,fill=\mausubsec,minimum width=1.5cm] (muccon) at (0,0){\bfseries A};
\node[text=white,left=3pt] at (muccon.east) {\bfseries #1};
\end{tikzpicture}%
}
%
\titlespacing{\subsection}{0cm}{0cm}{-5mm}[0cm]
\titleformat{\subsection}
{\normalfont\fontsize{12pt}{12pt}\fontfamily{qag}\selectfont\bfseries\color{white}}
{}
{0mm}
{%
\begin{tikzpicture}[line join=round,line cap=round]
%-- Nội dung chính --	
\node[inner sep=4pt,fill=\mausubsec,text=white,text width=\textwidth-8pt,align=center] (subsec) at (0,0){\MakeUppercase{\thesubsection. ~#1}};
\draw[line width=1pt,\mausubsec] ([yshift=-2pt]subsec.south west)--++(\textwidth,0);
\end{tikzpicture}
}
[]
%----------ĐN subsubsection-----------------------
\newcommand{\dnsubsub}[1]{%
\begin{tikzpicture}[baseline=(muccon.base),line join=round,line cap=round]%
\node[rectangle,inner sep=7pt,text=\mausubsubsec,fill=\mausubsubsec,minimum width=1.2cm] (muccon) at (0,0){\bfseries A};
\foreach \ii/\mau in {0/\mausubsubsec!30,0.3/\mausubsubsec!70,0.7/\mausubsubsec}{
\fill[\mau,draw=white,line width=3pt] (muccon.north west)--([xshift=0cm+1cm-\ii cm]muccon.north east)
--([xshift=1cm-\ii cm]muccon.south east)--(muccon.south west) --cycle;
}
\node[text=white,left=3pt] at (muccon.east) {\bfseries #1};
\end{tikzpicture}%
}
\titlespacing{\subsubsection}{0pt}{0mm}{0mm}[0cm]
\titleformat{\subsubsection}
{\fontsize{13pt}{18pt}\fontfamily{put}\selectfont\bfseries\color{\mausubsubsec}}
{\hspace*{-1mm}\dnsubsub{\Large\thesubsubsection}}
{2mm}
{#1}
[]
%----------ĐN paragraph-----------------------
\newcommand{\dnparagraph}[1]{
\begin{tikzpicture}[baseline=(muccon.base),line join=round,line cap=round]%
\node[rectangle,inner sep=3.5pt,text=\mausubsubsec,fill=\mausubsubsec,minimum width=1cm,rounded corners=5pt] (muccon) at (0,0){\bfseries 1};
\node[text=white] at (muccon.center) {\bfseries #1};
\end{tikzpicture}%
}
\titlespacing{\paragraph}{0pt}{0mm}{0mm}[0cm]
\titleformat{\paragraph}
{\fontsize{11.5pt}{17pt}\fontfamily{put}\selectfont\bfseries\color{\mausubsubsec}}
{\dnparagraph{\theparagraph.}}
{0.25em}
{#1}
[]
%============================
\def\itemKN{\color{\mauitemKN}\faCheckSquareO}
\def\itemCI{\color{\mauitemCI}\faCheckCircleO}
%%======= Thiết lập labelitem, labelenumerate
%\renewcommand{\labelitemi}{\color{red}\faCheckSquareO}
\renewcommand{\labelitemi}{\color{\mauitem}\faCheckCircleO}
\renewcommand{\labelitemii}{\color{\mauitem}\bf ---}
\renewcommand{\labelitemiii}{\color{\mauitem}\bf +}
\renewcommand{\labelenumi}{\alph{enumi})}
%\renewcommand{\labelenumii}{\color{blue}\bf\arabic{enumi}.\arabic{enumii}}
%============================
%============================
% Canh chỉnh mục lục chính
\setcounter{secnumdepth}{4} %Độ sâu đánh số
\setcounter{tocdepth}{2} %Độ sâu mục lục
\contentsmargin{0cm}
%~~~~~~~~~~~~~~~~~~~~~
\renewcommand*\l@part[2]{%
\ifnum \c@tocdepth >-2\relax
\addpenalty{-\@highpenalty}%
\addvspace{20pt \@plus\p@}%
\setlength\@tempdima{3em}%
\begingroup
\hypersetup{linkcolor=violet}
\tikz[remember picture, overlay]{
\fill[\mauPHAN] (0,0) rectangle +(\textwidth,0.85);
\draw (0,0.425) node[right=5pt]
{\color{white}\fontsize{13pt}{1pt}\fontfamily{qag}\selectfont\bfseries  {\scshape Phần} #1};
}
\par\smallskip
\penalty\@highpenalty
\endgroup
\fi
}
%------------------------
%\titlecontents{part}[0pc]
%{\addvspace{10pt}%
%	\color{red!70!black}\fontsize{18pt}{1pt}\fontfamily{qag}\selectfont\bfseries 
%}%
%{}
%{}
%{}%
%~~~~~~~~~~~~~~~~~~~~~
\titlecontents{chapter}[6.5pc] % nd cách trái
{\addvspace{0pt}%
\color{\mauCHUONG}\fontsize{13pt}{16pt}\fontfamily{put}\selectfont\bfseries
}%
{\contentslabel[\chaptertitlename\,\thecontentslabel.]{6.5pc}} %nhãn
{}
{\hfill\bfseries\thecontentspage
}%
%~~~~~~~~~~~~~~~~~~~~~
\titlecontents{section}[10pc]
{\addvspace{0pt}\bfseries\color{\mauSEC}}
{\fontsize{12.5pt}{15pt}\selectfont\sffamily\contentslabel[{Bài\,\thecontentslabel.}]{3.5pc}}
{}
{\hfill
\thecontentspage
}
[]
%~~~~~~~~~~~~~~~~~~~~~
\titlecontents{subsection}[10pc]
{\addvspace{0pt}\color{\mauSUBSEC}}
{\fontsize{12pt}{15pt}\selectfont\sffamily\contentslabel[\tron{\thecontentslabel}]{2.8pc}}
{}
{{\tiny\dotfill}\thecontentspage}
[]
%~~~~~~~~~~~~~~~~~~~~~
%--------------------------
% ĐỊNH NGHĨA CÁC MÔI TRƯỜNG 
%--------------------------
\listenumerate{dn,dl,tc,nx,ex}%xuống dòng khi liệt kê
\theoremstyle{plain} %
\theoremheaderfont{\scshape} %đầu
\theorembodyfont{\normalfont} % thân
\theoremseparator {.} % Ngăn cách
\newtheorem{mydn}{\color{\maudn}\faBolt\, Định nghĩa}%[section]
%===================================ĐNghĩa
\theoremstyle{plain} %
\theoremheaderfont{\fontfamily{put}\bfseries} %đầu
\theorembodyfont{\normalfont} % thân
\theoremseparator {.} % Ngăn cách
\newtheorem{vd}{\color{\mauVD}\damEX%\fontfamily{qag}\selectfont\bfseries
%\faToggleOn\ 
%\faUnlink\ 
\iconVD	Ví dụ}%[section]
\newtheorem{bt}{\color{\mauBT}%\fontfamily{qag}\selectfont\bfseries%
\damEX\iconVD Bài}
\newtheorem{EXp}{\color{\mauEX}%\fontfamily{qag}\selectfont\bfseries%
\damEX Câu}
%===================================
\theoremstyle{plain} %
\theoremheaderfont{\scshape} %đầu
\theorembodyfont{\slshape} % thân 
\theoremseparator {.} % Ngăn cách 
\newtheorem{mydl}{\color{\maudl}\faBolt\, Định lí}%[section]
\newtheorem{mytc}{\color{\mauhq!70!black}\faBolt\, Tính chât}%[section]
\newtheorem{myhq}{\color{\mauhq!70!black}\faBolt\, Hệ quả}%[section]
%====================================
\theoremstyle{nonumberplain} %ko đánh số, ko xuống dòng
\theoremheaderfont{\scshape} %đầu
\theorembodyfont{\normalfont} %phần thân
\theoremseparator {.} %ngăn cách
\newtheorem{mynx}{\color{\mauhq!70!black}\faBolt\, Nhận xét}
%====================================
%\theoremstyle{nonumberplain} %ko đánh số, ko xuống dòng
%\theoremheaderfont{\it} %đầu
%\theorembodyfont{\normalfont} %phần thân
%\theoremseparator {.} %ngăn cách
%\newtheorem{proof}{Chứng minh}
%\newenvironment{proof}{\noindent\textit{Chứng minh}. }{}
%====================================
\newenvironment{tomtat}{\noindent}{}
%====================================Hộp
\newenvironment{boxdl}
{\begin{tcolorbox}
[enhanced jigsaw,breakable,pad at break*=1mm,
colback=\maudl!10,boxrule=0pt,frame hidden, 
opacityback=0.5,
left=3mm, right=0pt, bottom=2pt, top=2pt,
before skip=2mm,
after skip=2mm,arc=0pt,
borderline west={1mm}{0mm}{\maudl!30}
%		overlay={\draw[line width=4pt,\maudl!50] ([xshift=1.5pt]interior.north west)--([xshift=1.5pt]interior.south west);}
]}
{\end{tcolorbox}}
\newenvironment{boxdn}
{\begin{tcolorbox}
[enhanced jigsaw,breakable,pad at break*=1mm,
%		drop fuzzy shadow=orange,
colback=\maudn!10,boxrule=0pt,frame hidden, 
opacityback=0.5,
left=3mm, right=0pt, bottom=2pt, top=2pt,
before skip=2mm,
after skip=2mm,arc=0pt,
borderline west={1mm}{0cm}{\maudn!30}
%		overlay={\draw[line width=4pt,\maudn!50] ([xshift=2pt]interior.north west)--([xshift=1.5pt]interior.south west);}
]}
{\end{tcolorbox}}

\newenvironment{boxkn}
{\begin{tcolorbox}
[enhanced jigsaw,breakable,pad at break*=1mm,
%		drop fuzzy shadow=orange,
colback=\mauhq!10,boxrule=0pt,frame hidden, 
opacityback=0.5,
left=3mm, right=0pt, bottom=2pt, top=2pt,
before skip=2mm,
after skip=2mm,arc=0pt,
borderline west={1mm}{0cm}{\mauhq!50}
%		overlay={\draw[line width=4pt,\mauhq!50] ([xshift=2pt]interior.north west)--([xshift=1.5pt]interior.south west);}
]}
{\end{tcolorbox}}
%--------------------Đn lại môi trường
\newcommand{\khungMT}[2]{%
\begin{tikzpicture}[baseline=(Akhung.base),scale=1]%
\node[fill=#1,rounded corners=4pt,inner sep=2pt] (Akhung) {\noindent\color{white}\fontfamily{put}\selectfont\bfseries\,#2\,};
\end{tikzpicture}%
}
%---
\renewenvironment{mydn}{
\stepcounter{mydn}\noindent\khungMT{\maudn}{\iconVD Khái niệm}%
}{}
\renewenvironment{mydl}{
\stepcounter{mydl}\noindent\khungMT{\maudl}{\iconVD Định luật}%
}{}
\renewenvironment{mytc}{
\stepcounter{mytc}\noindent\khungMT{\mauhq!70!black}{\iconVD Tính chất}%
}{}
\renewenvironment{myhq}{
\stepcounter{myhq}\noindent\khungMT{\mauhq!70!black}{\iconVD Hệ quả \themyhq}%
}{}
\renewenvironment{mynx}{
\stepcounter{mynx}\noindent\khungMT{\mauhq!70!black}{\iconVD Nhận xét \themynx}%
}{}
%---
\newenvironment{dn}{\begin{boxdn}\begin{mydn}}{\end{mydn}\end{boxdn}}
\newenvironment{dl}{\begin{boxdl}\begin{mydl}}{\end{mydl}\end{boxdl}}
\newenvironment{tc}{\begin{boxkn}\begin{mytc}}{\end{mytc}\end{boxkn}}
\newenvironment{nx}{\begin{boxkn}\begin{mynx}}{\end{mynx}\end{boxkn}}
\newenvironment{hq}{\begin{boxkn}\begin{myhq}}{\end{myhq}\end{boxkn}}
%--------------------Chú ý
\newenvironment{note}
{\begin{tcolorbox}
[enhanced jigsaw,breakable,pad at break*=1mm,
opacityback=0,boxrule=0pt,frame hidden,
left=8mm, right=0pt, bottom=0pt, top=0pt,
before skip=1mm,
after skip=1mm,
underlay unbroken and first={
\draw ([xshift=0.3cm,yshift=-0.32cm]interior.north west) node[\mauly]{\large\bfseries \faExclamationTriangle};
},
fontupper=\it,
]}
{\end{tcolorbox}}
%\let\mynote\note
%\renewcommand{\note}{\mynote{\bfseries\color{\mauly}Lưu ý:}} 
%--------------------ghi chú
\newenvironment{ghichu}
{\begin{tcolorbox}
[enhanced jigsaw,breakable,pad at break*=1mm,
opacityback=0,boxrule=0pt,frame hidden,
left=8mm, right=0pt, bottom=0pt, top=0pt,
before skip=2mm,
after skip=1mm,
underlay unbroken and first={
\draw ([xshift=0.3cm,yshift=-0.32cm]interior.north west) node[\maubl,yscale=-1]{\large\bfseries \faPaperclip};
},
%		fontupper=\it,
]}
{\end{tcolorbox}}
\let\myghichu\ghichu
\renewcommand{\ghichu}{\myghichu{\bfseries\color{\maubl}Ghi chú:}}
%--------------------Chú ý
\newenvironment{luuy}
{\begin{tcolorbox}
[enhanced jigsaw,breakable,pad at break*=1mm,
opacityback=0,boxrule=0pt,frame hidden,
left=8mm, right=0pt, bottom=0pt, top=0pt,
before skip=1mm,
after skip=1mm,
underlay unbroken and first={
\draw ([xshift=0.3cm,yshift=-0.32cm]interior.north west) node[\mauly]{\large\bfseries \faExclamationTriangle};
},
fontupper=\it,
]}
{\end{tcolorbox}}
%--------------------Bình luận
\newenvironment{binhluan}
{\begin{tcolorbox}
[enhanced jigsaw,breakable,pad at break*=1mm,
opacityback=0,boxrule=0pt,frame hidden,
left=6mm, right=0pt, bottom=0pt, top=0pt,
before skip=2mm,
after skip=2mm,
%		underlay unbroken and first={
%			\draw ([xshift=0.3cm,yshift=-0.25cm]interior.north west) node[magenta]{\large\bfseries 
	%				%\faBell
	%				\faEdit
	%			};
%		},
fontupper=\it,
]}
{\end{tcolorbox}}
\let\mybinhluan\binhluan
\renewcommand{\binhluan}{\mybinhluan{\hspace*{-0.75cm}\bfseries\color{\maubl}\faEdit\,\hspace*{0.15cm} Bình luận:}} 
%==================
%--------------------đáp số
\renewcommand{\dapso}[1]{\hspace{\fill}\mbox{}\linebreak[0]\hspace*{\fill}\mbox{\fontsize{8pt}{1pt}\selectfont\color{blue!80!black}{\color{red}{\faKey}}~#1}
}
%====================
\setlength{\parindent}{0pt} %không thụt đầu dòng
\usepackage[final]{pdfpages} % input file bìa pdf
%----Cài đặt \boxEX -ô vuông sd tài liệu
\def\hideboxEX{}
\renewcommand{\boxEX}[2][15pt]{%
\setbox1=\hbox{#2\strut}
\def\sizeboxEX{#1}%Thay #1 bằng 15pt để cố định tất cả boxEX
\ \EXbox{\hideboxEX#2\strut}\ %
}
%===Định nghĩa in Mục lục chính
\def\muclucchinh{
\pagenumbering{roman}%đánh số trang dạng i,ii,...
\thispagestyle{empty}
\tableofcontents %lệnh in mục lục chính
%	\cleardoublepage %tạo trang trống
}
%--------------Cài đặt lại dòng kẻ \dotline
%--Môi trường thủ công 2 cột Dotted
\newcommand{\dongcham}[1]{
\def\sod{#1}
\pgfmathsetmacro{\sodong}{2*\sod -1} 
\columnsep=10pt
\vspace*{-3.5mm}
\begin{multicols}{2}
\foreach \dotline in{1,...,\sodong}
{\noindent\color{\maucham}{\tiny\dotfill}\\[1mm]}
{\noindent\color{\maucham}{\tiny\dotfill}}\\[-4mm]
\end{multicols}
}
%---ĐN dòng chấm kiểu 1 cột
\def\dotlineOne{
\pgfmathsetmacro{\sodong}{\numlinedot-1} 
\par
\foreach \dotline in{1,...,\sodong}
{\noindent\color{\maucham}{\tiny\dotfill}\\[1mm]}
{\noindent\color{\maucham}{\tiny\dotfill}}\\[-2mm]
}
%--Đn dòng chấm kiểu 2 cột
\def\dotlineTwo{
\pgfmathsetmacro{\sodong}{2*\numlinedot+1} 
\columnsep=5pt
%		\setlength{\columnseprule}{0.4pt} % đường kẻ giữa
%		\renewcommand{\columnseprulecolor}{\color{\maucham}} %màu đường kẻ
\vspace*{-3.5mm}
\begin{multicols}{2}
\foreach \dotline in{1,...,\sodong}
{\noindent\color{\maucham}\tiny\dotfill\\[1mm]}
{\noindent\color{\maucham}\tiny\dotfill}\\[-4mm]
\end{multicols}
}
%--SD 2 kiểu đn trên vào lệnh dotlineEX
\renewcommand{\dotlineEX}[1]{
\setlength{\parindent}{0pt}%
\def\numlinedot{#1}
\dotlineTwo % CHỌN KIỂU DÒNG CHẤM
%	\dotlineOne
}
%--SD trong foot head
\newcommand{\dotlineEXhead}[1]{
\setlength{\parindent}{0pt}%
\def\numlinedot{#1}
\dotlineOne % CHỌN KIỂU DÒNG CHẤM
}
%--
\renewcommand{\dotlineans}[2]{
\hideansEX{#2}
\AfterEndEnvironment{#2}{%
\dotlineEX{#1}
}}
%---------ĐN các khung câu hỏi
%%-- ẩn / hiện mã QR đề thi
\def\maQR{}
\newcommand{\hienmaQR}{
\def\maQR{\qrcode[hyperlink,height=0.8cm,version=1]{\linkdethi}}
}
\newcommand{\anmaQR}{}
%Lệnh chung
%KHUNG SỐ 1
\newcommand{\khungcauhoiMOT}[4]{
\begin{tcolorbox}
[
enhanced jigsaw,breakable,pad at break*=1mm,
%	frame hidden, opacityback=0,
colback=#2!6,
colframe=#2,
drop fuzzy shadow=#2!30!gray, %đổ bóng
opacityback=0.5,
left=2mm,right=2mm,top=2pt,bottom=2mm,
%--
before skip=2mm,
after skip=3mm,
boxrule=0.5pt, %đường viền dày
%	toprule=1pt,
%bottomrule=1pt,
%		leftrule=2.5pt,
%rightrule=0pt,
arc=0mm,
fonttitle=\fontfamily{qag}\bfseries\selectfont,
title={\damEX\iconVD #1 \stepcounter{#3}#4},
attach boxed title to top left={xshift=0cm,yshift=-\tcboxedtitleheight},
boxed title style={empty},
underlay boxed title={
\fill[#2] (title.north west)--(title.south west)--(title.south east) 
to[out=5,in=-175] ([xshift=0.5cm]title.north east)
-- cycle;
\draw[#2] ([yshift=-2pt]title.south west)--([xshift=0.5pt, yshift=-2pt]title.south east) 
to[out=5,in=-175] ([xshift=0.5cm+2.5pt]title.north east)
;
},
underlay unbroken and first={
\node[anchor=north east,inner sep=0pt,#2] at (frame.north east) {
	\maQR
};
}	
]}
%KHUNG SỐ 2
\newcommand{\khungcauhoiHAI}[4]{
\begin{tcolorbox}
[
enhanced jigsaw,breakable,pad at break*=1mm,
breakable,  
frame hidden, 
%	frame hidden, opacityback=0,
colback=#2!3,
colframe=#2,
drop fuzzy shadow=#2!30!gray, %đổ bóng
opacityback=0.5,
left=2mm,right=2mm,top=2pt,bottom=2mm,
%--
before skip=2mm,
after skip=3mm,
boxrule=0pt, %đường viền dày
%	toprule=1pt,
%bottomrule=1pt,
%		leftrule=2.5pt,
%rightrule=0pt,
arc=0mm,
%		fonttitle=\fontfamily{qag}\bfseries\selectfont,
title={\damEX\iconVD #1\,\, \stepcounter{#3}#4},
attach boxed title to top left={xshift=3mm,yshift=-\tcboxedtitleheight},
boxed title style={%
	enhanced, 
	colframe=#2, 
	colback=#2, 
	borderline east={1.5pt}{3.5pt}{white},
	borderline east={3.5pt}{0pt}{white},    
	borderline east={2.75pt}{0.25pt}{#2!60},    
	borderline east={2.75pt}{-4pt}{#2!30},    
	bottomrule=0pt, 
	rightrule=5pt, 
	sharp corners, 
}, 
%		interior style={}, 
underlay boxed title={
	\draw[#2, line width=2pt] (title.west)++(0,2pt)--++(-6mm,0)
	(title.west)++(0,-2pt)--++(-6mm,0);
},
overlay={%
	\begin{scope}[line width=2pt, #2]
		\draw
		(frame.north west) coordinate (iniPoint)
		-- (iniPoint|-frame.south west);
\end{scope}},
underlay unbroken and first={
	\node[anchor=north east,inner sep=0pt,#2] at (frame.north east) {
		\maQR
	};
}
]}
%BỘ ĐẾM PHỤ
\newcounter{EX}\setcounter{EX}{0}\newcounter{BT}\setcounter{BT}{0}\newcounter{VD}\setcounter{VD}{0}
%%%%%%%% KIỂU TUỲ CHỌN KHUNG CÂU HỎI
\newcommand{\TuyChonKhung}[1]{
\def\chonkhung{#1}
\def\mot{1}
\def\hai{2}
\ifx\chonkhung\mot%
\def\beginboxEX{\khungcauhoiMOT{CÂU}{\mauEX}{EX}{\theEX}}
\def\beginboxBT{\khungcauhoiMOT{BÀI}{\mauBT}{BT}{\theBT}}
\def\beginboxVD{\khungcauhoiMOT{VÍ DỤ}{\mauVD}{VD}{\theVD}}
\else
\def\beginboxEX{\khungcauhoiHAI{CÂU}{\mauEX}{EX}{\theEX}}
\def\beginboxBT{\khungcauhoiHAI{BÀI}{\mauBT}{BT}{\theBT}}
\def\beginboxVD{\khungcauhoiHAI{VÍ DỤ}{\mauVD}{VD}{\theVD}}
\fi
}
%---------Khoảng cách đoạn
\setlength{\parskip}{2mm}
%--------------------------Đn lại màu sắc theo từng môi trường
%\AtBeginEnvironment{bt}{\def\mauLG{\mauBT} \def\mauDA{\mauBT} \def\mauTrue{\mauBT} \def\maucham{\mauBT}}
%\AtBeginEnvironment{ex}{\def\mauLG{\mauEX} \def\mauDA{\mauEX} \def\mauTrue{\mauEX} \def\maucham{\mauEX}}
%\AtBeginEnvironment{vd}{\def\mauLG{\mauVD} \def\mauDA{\mauVD} \def\mauTrue{\mauVD} \def\maucham{\mauVD}}
%===========================MainDeThi
%FULL WIDTH
\def\FULLWIDTH{
\newpage
\fancyhead[LO,RE]{\headmucluc}
\def\headchuong{\headmucluc}
\newgeometry{top=\tren cm, bottom=\duoi cm, left=\trai cm, right=\phai cm}
\def\mepphai{\phai}
}
\def\NOTE{
\newpage
\fancyhead[LO,RE]{\myfancyhead}
\def\headchuong{\myfancyhead}
\pgfmathsetmacro{\mepphai}{\phai+\leftnote}
\newgeometry{top=\tren cm, bottom=\duoi cm, left=\trai cm, right=\mepphai cm}
}
%==================
%=======Đn các phương án ABCD
\def\khoanhtrondapan{
\renewcommand*\circled[1]{\tikz[baseline=(char.base)]{
		\node[shape=circle,draw=\mauDA,inner sep=1pt] (char) {##1};}}
%	\renewcommand{\TrueEX}{\stepcounter{dapan}
	%		{\squareEX{\textbf{\damEX\color{\mauDA}\Alph{dapan}}}} \ignorespaces}
\renewcommand{\FalseEX}{\stepcounter{dapan}
	{\circled{\textbf{\damEX\color{\mauDA}\Alph{dapan}}}} \ignorespaces}
	}
	%----
	\def\totrondapan{
\renewcommand*\circled[1]{\tikz[baseline=(char.base)]{
		\node[shape=circle,draw=\mauDA,fill=\mauDA,inner sep=1pt] (char) {##1};}}
%	\renewcommand{\TrueEX}{\stepcounter{dapan}
	%		{\squareEX{\textbf{\damEX\color{\mauDA}\Alph{dapan}}}} \ignorespaces}
\renewcommand{\FalseEX}{\stepcounter{dapan}
	{\circled{\textbf{\damEX\color{white}\Alph{dapan}}}} \ignorespaces}
	}
	%----
	\def\khongkhoanhtrondapan{
%	\renewcommand{\TrueEX}{\stepcounter{dapan}
	%		{\squareEX{\textbf{\damEX\color{\mauDA}\Alph{dapan}}}} \ignorespaces}
\renewcommand{\FalseEX}{\stepcounter{dapan}
	{\textbf{\damEX\color{\mauDA}\Alph{dapan}.}} \ignorespaces}
	}
	%%%%%%%%%%%%%%%%%% Khung tiêu đề, khung đề thi, đề cương
	%Tạo mã qrcode
	%----------gói tạo mã qr
	\usepackage{qrcode}
	%----------
	\def\linkdecuong{https://toanvip307.blogspot.com/}
	\def\linkdethi{https://toanvip307.blogspot.com/}
	%---
	%\def\mauname{teal}\def\mauSO{red}
	\newcounter{decuongso}\setcounter{decuongso}{0}
	\font\dam=ugqb8v at 15pt
	\font\damTT=ugqb8v at 17pt
	\renewcommand{\hrulefill}{%
\leavevmode\leaders\hrule height 1pt\hfill\kern0pt }
%---Đn name
\newcommand{\decuong}[4]{%
\stepcounter{decuongso}%
\setcounter{ex}{0}\setcounter{bt}{0}\setcounter{vd}{0}%
\boldmath\fontfamily{qag}\selectfont\color{\mauname}%
%%--trước khung decuong
{\fontsize{10}{11}\selectfont\raggedleft 
	\parbox{3.5cm}{Ngày: \hrulefill}
	\hspace*{1mm}
	\parbox{0.75\textwidth-4cm}{Họ và tên: \hrulefill}
	\hspace*{1mm}
	\parbox{0.25\textwidth}{Lớp: \hrulefill}
}\\[-2mm]
%%--khung decuong
\begin{tcolorbox}[boxrule=0.7pt,arc=0mm,breakable,colframe=\mauSO,colback=\mauname!2,before skip=0mm,after skip=2mm, left=0mm,right=0mm,top=2mm,bottom=2mm]\color{\mauname}
	\begin{minipage}[c]{3.5cm}
		\begin{center}
			\qrcode[hyperlink,height=2.7cm,version=2]{\linkdecuong}%\\
			%				{\fontsize{8}{0}\selectfont\sf Phần hướng dẫn giải}
		\end{center}
	\end{minipage}
	\hspace{0.2cm}
	\begin{minipage}[c]{\textwidth-4.4cm}
		\begin{center}
			%%---
			{\damTT\color{\mauSO} \MakeUppercase{#1}}\\[2pt]
			{\color{\mauSO}\fontsize{14}{11}\selectfont\bfseries \MakeUppercase{#2}}\\[2pt]
			{\color{\mauSO}\fontsize{11}{11}\selectfont\bfseries #3}\\[2pt]
			{\fontsize{14}{11}\selectfont\bfseries \MakeUppercase{#4}}\\[-1mm]
		\end{center}
	\end{minipage}
\end{tcolorbox}
%%--- Phần note đầu đề
%	\notename
\vspace*{0.5cm}
\addcontentsline{toc}{section}{\hspace*{-4.2cm}\sf Phiếu số \thedecuongso: #4} % đưa MT vào mục lục
}
%--Sang trang 
\BeforeBeginEnvironment{decuong}{
\ifnum\the\value{decuongso}>0
\newpage
\fi
}
%%%%%%%%%%%%%%%%%%% Khung đề thi
\newcounter{deso}\setcounter{deso}{0}
\font\damDETHI=ugqb8v at 13pt
\newcommand{\dethi}[3]{%
\setcounter{page}{1}%
\setcounter{ex}{0}\setcounter{bt}{0}\setcounter{vd}{0}%
\setcounter{EX}{0}\setcounter{BT}{0}\setcounter{VD}{0}%
\stepcounter{deso}
\label{\thedeso}
\begin{minipage}[t]{0.5\textwidth}
	\centering
	\vspace*{1mm}
	{\color{\mauSO}\damDETHI PHÒNG GD\&ĐT\ \MakeUppercase{#1}}\\	
	{\fontsize{12}{12}\color{\mauname}\bfseries\fontfamily{qag}\selectfont \boxed{\textbf{#2}}}\\
	\fontsize{12}{12}\color{\mauname}\fontfamily{qag}\selectfont{\it (Đề thi gồm có 0\pageref{de\thedeso} trang)}
\end{minipage}
%%---
\begin{minipage}[t]{0.5\textwidth}
	\centering
	{\color{\mauname}\damDETHI \MakeUppercase{#3}}\\
	{\color{\mauname}\fontfamily{qag}\fontsize{12}{12}\selectfont\bfseries Năm học: 2023 - 2024}\\
	{\color{\mauname}\fontfamily{qag}\fontsize{12}{12}\selectfont\bfseries Môn: Toán}\\
	{\color{\mauname}\fontfamily{qag}\fontsize{11}{12}\selectfont Thời gian: \textit{90 phút} (không kể phát đề)}\\
	%		\rule[2.2ex]{1.5in}{1pt}
\end{minipage}
\\
\noindent\fontsize{11}{12}\fontfamily{qag}\selectfont\color{\mauname}
\parbox[t]{\textwidth-2cm-5cm-7pt}{
	\text{Họ và tên thí sinh: }{\tiny\dotfill}
}
\parbox[t]{5cm}{
	\text{Số báo danh: }{\tiny\dotfill}
}
\parbox[b]{2cm}{
	\setlength\fboxrule{2pt}
	\setlength\fboxsep{5pt}
	\hfill\vspace*{-2mm}\fbox{\fontfamily{qag}\selectfont\bfseries Đề số \thedeso}
}
%-- gạch dưới
\centerline{
	\pgfornament[width=4cm,color=\mauname]{88}
}\\[1mm]
%%---
\addcontentsline{toc}{section}{\hspace*{-4.2cm}\sf Đề số \thedeso: #3 môn Toán, #1} % đưa MT vào mục lục
}
%--Sang trang 
\BeforeBeginEnvironment{dethi}{
\ifnum\the\value{deso}>0
\newpage
\fi
}
%%%%%%%%%%%%%%%% Thầy Live Hoá
\newcommand{\loigiaiTT}[2]{
\begin{minipage}[t]{0.4\textwidth-2mm}
	\begin{khung4}{Gợi ý}
		#1
	\end{khung4}
\end{minipage}\hspace*{2mm}
\begin{minipage}[t]{0.6\textwidth}\vspace*{5mm}
	#2
\end{minipage}	
}
%%=============ĐN HIỆN CÂU EX CẦN THIẾT if
%\newcommand{\hienEXS}[2]{
%	\def\biendau{#1}\def\biencuoi{#2}%
%	\pgfmathsetmacro{\sodau}{int(\biendau-1)}
%	\pgfmathsetmacro{\socuoi}{int(\biencuoi+1)}
%	\setcounter{EXp}{#1-1}
%	\RenewEnviron{ex}{
	%		\stepcounter{ex}%
	%		\ifnum\value{ex}<\socuoi
	%		\ifnum\value{ex}>\sodau
	%		\par%
	%		\begin{EXp}
		%			\BODY% 
		%		\end{EXp}
	%		\fi\fi
	%	}
%	%
%	\AtEndEnvironment{dethi}{\setcounter{EXp}{#1-1}}
%	%
%	\AtEndEnvironment{EXp}{
	%		\ifnum\the\value{numTrue}=1
	%		\scantokens{\begin{EXsol}A\end{EXsol}}
	%		\fi
	%		\ifnum\the\value{numTrue}=2
	%		\scantokens{\begin{EXsol}B\end{EXsol}}
	%		\fi
	%		\ifnum\the\value{numTrue}=3
	%		\scantokens{\begin{EXsol}C\end{EXsol}}
	%		\fi
	%		\ifnum\the\value{numTrue}=4
	%		\scantokens{\begin{EXsol}D\end{EXsol}}
	%		\fi
	%		\setcounter{numTrue}{0}
	%	}
%}
%%%%%%%%%%%%%%%%%%%%%%%
\newcolumntype{M}[1]{>{\centering\arraybackslash}m{#1}} % m là middle
\newcolumntype{C}[1]{>{\centering\arraybackslash}p{#1}} 
\newcolumntype{L}[1]{>{\raggedright\arraybackslash}p{#1}}
% lệnh siunit 
\newcommand{\xsi}[2]{\SI[parse-numbers=false]{#1}{#2}}